\backhead{FURTHER EXERCISES}{khred}
\begin{exlist}
\item Using the rate-competition law, find the $K_s$ at which the memory fate of drainage is 0.1, 0.5 and 0.9.\exans{7.9, 71.5, 643 cm\,d$^{-1}$}
\item The memory fates of a soil are $\phi_D=0.29$, $\phi_T=0.58$, $\phi_E=0.04$. What is the remainder, and what does it represent?\exans{0.09; runoff and storage change}
\item A figure on a log axis shows points at 10, 100 and 1000. Are they equally spaced? What ratio does one decade represent?\exans{yes; a factor of 10}
\item Compute the RMSE and the bias of the predictions 0.30, 0.55, 0.12 against observations 0.25, 0.60, 0.10.\exans{0.042; $+0.007$}
\item A recoverable fraction is $1-\psi_D$. Find it for $\psi_D=0.31$ and 0.71, and the recovered water of a saving of 120 mm.\exans{0.69 and 0.29; 82.8 and 34.8 mm}
\end{exlist}
\backhead{ESSAY QUESTIONS}{khred}
{\small\itshape Answer each question in about one page, with equations, a sketch and at least one number from this chapter.}\par\vspace{3pt}
\begin{exlist}
\item Explain how to read a figure of the atlas critically, using one figure of this chapter as an example.
\item Discuss the rate-competition law: its form, its evidence and its limits.
\item Explain the difference between a memory fate and a loss fate.
\item Discuss what the coupled fluxes of the study show about the closure of the water balance.
\item Explain why the recoverable fraction differs between soils.
\item Choose one figure of this chapter and write the paragraph of a paper that describes it.
\end{exlist}
